\section{Consequences and open questions}\label{sec:discussion}

Appendix~\ref{sec:counts} gives unconditional logical upper bounds for the total count: a finite upper cut has a $\forall n\,\exists m$ description with a decidable finite test. Under $C_T$, Theorem~\ref{thm:completeness} supplies matching $\Pi^0_2$ and $\Sigma^0_2$ completeness at the rate $\eps_1$. The stationary diagonal-coupling argument also gives one-sided second-level hardness and excludes nontrivial $\Delta^0_2$ predicates of $\sigma_{\mathrm{stat}}$. The total-count cuts are the concrete complete instances. Arithmetic upper bounds for the ordered and Bragg subset counts and $L$, and matching degrees for all-law ergodicity, remain open.

The non-tiling correspondence preserves each corresponding state's Bragg spectrum. On the tiling branch, however, the universal crystalline count records a peak somewhere; the construction forces the period-two wavevector $q_\pi$. Counts at other specified wavevectors can retain host information and are outside the Rice assertion. A Rice theorem for the full stationary convex set or its translation dynamics would likewise require a tiling branch that removes more host information.

Realizability of exactly $k\ge2$ extreme TI stationary states, smaller fixed alphabets, and lower dimensions remain open. The totality reduction uses a stationary coupling property; the stronger all-input forgetting needed to retain $E$ is a separate issue.

